% GENERATED FILE. Edit canonical lessons, then run npm run generate:books.
% canonical-source-hash: fnv1a64:e0c3655ada634104
% canonical-lessons: ES-C664-sabana, ES-C664-balcon, ES-C664-vista, ES-C664-enchufe, ES-C664-grifo

\chapter{Cómo es la habitación}
\label{ch:es-como-es-la-habitacion}
\hlchaptermodality{664}

\begin{chapteropening}
\textbf{By the end of this chapter:} \emph{I can describe a hotel room: the sheets, the balcony, the view, the sockets and the taps.}

Complete the last lesson of chapter 664: I can describe a hotel room: the sheets, the balcony, the view, the sockets and the taps.
\end{chapteropening}

\section[la sábana]{la sábana — the sheet (on a bed)}

\label{lesson:ES-C664-sabana}



\begin{quote}
Before the new one: say the Spanish for the refund, then the Spanish for to refund.
\end{quote}

\subsection*{You'll want to know: la sábana}
\textbf{la sábana} — \textquotedblleft{}the sheet (on a bed)\textquotedblright{}.

\begin{quote}
Cambian las sábanas cada día.
\end{quote}

\begin{quote}
\emph{They change the sheets every day.}
\end{quote}

Do not mix it up with \emph{la sabana}, savannah, with no written accent.

\begin{cousinweb}
\emph{sábana} is inherited from Latin \emph{sabana}, the plural of \emph{sabanum} \textquotedblleft{}a linen cloth, a towel\textquotedblright{} (Lewis \& Short; EtymDB), from Greek \emph{sábanon}.
\end{cousinweb}

\begin{grammarlens}[title={What you've built}]
One more word for telling, explaining and describing at B1.
\end{grammarlens}

\subsection*{Your turn}
\begin{itemize}
\raggedright
  \item \emph{Say it:} \emph{la sábana}
  \item \emph{Say it:} \emph{la sábana}, once more
  \item \emph{Say it:} a sentence of your own with \emph{la sábana}
  \item \emph{Recall:} say the Spanish for the refund, then the Spanish for to refund, then say \emph{la sábana} again
\end{itemize}

\begin{tcolorbox}[breakable,colback=teal!4,colframe=teal!35!black,title={Before you move on}]
What does la sábana mean? (\textquotedblleft{}The sheet (on a bed)\textquotedblright{}.) Say it once more.
\end{tcolorbox}
\section[el balcón]{el balcón — the balcony}

\label{lesson:ES-C664-balcon}



\begin{quote}
Before the new one: say the Spanish for to refund, then the Spanish for the sheet.
\end{quote}

\subsection*{You'll want to know: el balcón}
\textbf{el balcón} — \textquotedblleft{}the balcony\textquotedblright{}.

\begin{quote}
La habitación tiene un balcón al mar.
\end{quote}

\begin{quote}
\emph{The room has a balcony over the sea.}
\end{quote}

\emph{Salir al balcón} — to go out on the balcony.

\begin{cousinweb}
\emph{balcón} was borrowed from Italian \emph{balcone}, which comes from Langobardic \emph{balko} \textquotedblleft{}beam\textquotedblright{}, a Germanic word (EtymDB).
\end{cousinweb}

\begin{grammarlens}[title={What you've built}]
One more word for telling, explaining and describing at B1.
\end{grammarlens}

\subsection*{Your turn}
\begin{itemize}
\raggedright
  \item \emph{Say it:} \emph{el balcón}
  \item \emph{Say it:} \emph{el balcón}, once more
  \item \emph{Say it:} a sentence of your own with \emph{el balcón}
  \item \emph{Recall:} say the Spanish for to refund, then the Spanish for the sheet, then say \emph{el balcón} again
\end{itemize}

\begin{tcolorbox}[breakable,colback=teal!4,colframe=teal!35!black,title={Before you move on}]
What does el balcón mean? (\textquotedblleft{}The balcony\textquotedblright{}.) Say it once more.
\end{tcolorbox}
\section[la vista]{la vista — the view; sight, eyesight}

\label{lesson:ES-C664-vista}



\begin{quote}
Before the new one: say the Spanish for the sheet, then the Spanish for the balcony.
\end{quote}

\subsection*{You'll want to know: la vista}
\textbf{la vista} — \textquotedblleft{}the view; sight, eyesight\textquotedblright{}.

\begin{quote}
Desde aquí hay una vista preciosa.
\end{quote}

\begin{quote}
\emph{There is a lovely view from here.}
\end{quote}

\emph{Con vistas al mar} — with a sea view. \emph{Hasta la vista} — see you.

\begin{cousinweb}
\emph{vista} is built on \emph{visto} \textquotedblleft{}seen\textquotedblright{}, the past form of \emph{ver} \textquotedblleft{}to see\textquotedblright{}, which comes from Latin \emph{vidēre} (EtymDB).
\end{cousinweb}

\begin{grammarlens}[title={What you've built}]
One more word for telling, explaining and describing at B1.
\end{grammarlens}

\subsection*{Your turn}
\begin{itemize}
\raggedright
  \item \emph{Say it:} \emph{la vista}
  \item \emph{Say it:} \emph{la vista}, once more
  \item \emph{Say it:} a sentence of your own with \emph{la vista}
  \item \emph{Recall:} say the Spanish for the sheet, then the Spanish for the balcony, then say \emph{la vista} again
\end{itemize}

\begin{tcolorbox}[breakable,colback=teal!4,colframe=teal!35!black,title={Before you move on}]
What does la vista mean? (\textquotedblleft{}The view; sight, eyesight\textquotedblright{}.) Say it once more.
\end{tcolorbox}
\section[el enchufe]{el enchufe — the socket, the plug; connections (influence)}

\label{lesson:ES-C664-enchufe}



\begin{quote}
Before the new one: say the Spanish for the balcony, then the Spanish for the view.
\end{quote}

\subsection*{You'll want to know: el enchufe}
\textbf{el enchufe} — \textquotedblleft{}the socket, the plug; connections (influence)\textquotedblright{}.

\begin{quote}
No hay enchufe al lado de la cama.
\end{quote}

\begin{quote}
\emph{There is no socket next to the bed.}
\end{quote}

Colloquially, \emph{tener enchufe} is to have connections that get you a job.

\begin{cousinweb}
\emph{enchufe} is built on \emph{enchufar} \textquotedblleft{}to fit one pipe into another\textquotedblright{} (EtymDB). The origin of the verb is not settled; it may be imitative.
\end{cousinweb}

\begin{grammarlens}[title={What you've built}]
One more word for telling, explaining and describing at B1.
\end{grammarlens}

\subsection*{Your turn}
\begin{itemize}
\raggedright
  \item \emph{Say it:} \emph{el enchufe}
  \item \emph{Say it:} \emph{el enchufe}, once more
  \item \emph{Say it:} a sentence of your own with \emph{el enchufe}
  \item \emph{Recall:} say the Spanish for the balcony, then the Spanish for the view, then say \emph{el enchufe} again
\end{itemize}

\begin{tcolorbox}[breakable,colback=teal!4,colframe=teal!35!black,title={Before you move on}]
What does el enchufe mean? (\textquotedblleft{}The socket, the plug; connections (influence)\textquotedblright{}.) Say it once more.
\end{tcolorbox}
\section[el grifo]{el grifo — the tap, the faucet}

\label{lesson:ES-C664-grifo}



\begin{quote}
Before the new one: say the Spanish for the view, then the Spanish for the socket.
\end{quote}

\subsection*{You'll want to know: el grifo}
\textbf{el grifo} — \textquotedblleft{}the tap, the faucet\textquotedblright{}.

\begin{quote}
El grifo del baño gotea.
\end{quote}

\begin{quote}
\emph{The bathroom tap is dripping.}
\end{quote}

Latin America says \emph{la llave} or \emph{la canilla}.

\begin{cousinweb}
\emph{grifo} is the same word as \emph{grifo} \textquotedblleft{}griffin\textquotedblright{}, from Latin \emph{gryphus} and Greek \emph{grýps} (EtymDB); old taps were often shaped like a griffin head.
\end{cousinweb}

\begin{grammarlens}[title={What you've built}]
One more word for telling, explaining and describing at B1.
\end{grammarlens}

\subsection*{Your turn}
\begin{itemize}
\raggedright
  \item \emph{Say it:} \emph{el grifo}
  \item \emph{Say it:} \emph{el grifo}, once more
  \item \emph{Say it:} a sentence of your own with \emph{el grifo}
  \item \emph{Recall:} say the Spanish for the view, then the Spanish for the socket, then say \emph{el grifo} again
\end{itemize}

\begin{tcolorbox}[breakable,colback=teal!4,colframe=teal!35!black,title={Before you move on}]
What does el grifo mean? (\textquotedblleft{}The tap, the faucet\textquotedblright{}.) Say it once more.
\end{tcolorbox}
